\documentclass{article}
\usepackage{import}
\import{../../template/}{article-template}
\graphicspath{{../../images/}}
\addbibresource{../../template/library.bib}

\title{Shell}
\author{PENG}
\date{\today}

\begin{document}
\maketitle
\tableofcontents
\section{Introduction}
\url{https://www.gnu.org/software/bash/manual/bash.html}

\section{Shell Syntax}
\subsection{Quoting}
\begin{definition}[quoting]
    Quoting is used to remove the special meaning of certain characters or words to the shell.
    \begin{verbatim}
echo -e "\t"
# \n : newline
# \t : horizontal tab
# \r : carriage return
\end{verbatim}
\end{definition}

\subsection{Pipelines}
\begin{definition}[pipeline]
    cmd1 output $\to$ cmd2 input
    \begin{verbatim}
cmd1 | cmd2      # stdout
cmd1 |& cmd2     # stderr to stdout
cmd1 2>&1 | cmd2 # stderr to stdout
\end{verbatim}
\end{definition}

\subsection{Lists of Commands}
\begin{definition}[list]~
    \begin{verbatim}
cmd1 && cmd2 # And, cmd1 true then cmd2
cmd1 || cmd2 # Or, cmd1 false then cmd2
cmd1; cmd2   # or terminated by newline
cmd1 & cmd2  # cmd1 async
\end{verbatim}
\end{definition}

\subsection{Construct}
\subsubsection{Looping}
\begin{definition}[looping]~
    \begin{verbatim}
for item in item1 item2 ...; do
    cmd
done
\end{verbatim}
\end{definition}

\subsubsection{Conditional}
\begin{definition}[conditional]~
    \begin{verbatim}
if [[ cond1 ]]; then
    cmd1
elif [[ cond2 ]]; then
    cmd2
else
    cmd3
fi

# cond:
# ==
# !=
# =~ : unquoted regex
\end{verbatim}
\end{definition}

\subsubsection{Grouping}
\begin{definition}[grouping]~
    \begin{verbatim}
( expr )  # subshell env, def var vanish when exit
{ expr; } # ; or newline is required
\end{verbatim}
\end{definition}

\subsection{Functions}
\begin{definition}[function]~
    \begin{verbatim}
func () {
    local var = $1
}
\end{verbatim}
\end{definition}

\subsection{Parameters}
\begin{definition}[parameters]~
    \begin{itemize}
        \item \verb|${N}| : position param
        \item \verb|$@| : expands to all pos params, \verb|$1 $2 ...|
        \item \verb|$0| : file or shell name
    \end{itemize}
\end{definition}

\subsection{Expansion}
\begin{definition}[brace expansion]
    Generate strings sharing a common prefix and suffix.
    \begin{verbatim}
echo file{1,2} # file1 file2
\end{verbatim}
\end{definition}

\begin{definition}[tidle expansion]~
    \begin{verbatim}
~ # $HOME
\end{verbatim}
\end{definition}

\begin{definition}[parameter expansion]~
    \begin{verbatim}
${param}

${param:-default}

${array[@]} # member of array

${!name[@]} # list of array indices
\end{verbatim}
\end{definition}

\begin{definition}[command substitution]~
    \begin{verbatim}
$(cmd)
`cmd`
\end{verbatim}
\end{definition}

\begin{definition}[arithmetic expansion]~
    \begin{verbatim}
$((expr)) # integer-only
\end{verbatim}
\end{definition}

\begin{definition}[filename expansion]~
    \begin{verbatim}
*  # any string
?  # single char
[] # char enclosed between []
*()
?()
+()
!(pattern list) # match except given patterns
\end{verbatim}
\end{definition}

\subsection{Redirection}
\begin{definition}[redirection]~
    \begin{verbatim}
ls > file 2>&1   # stdout and stderr to file
ls &> file       # or
pwd >> file 2>&1 # append stdout and stderr to file
pwd &>> file     # or
\end{verbatim}
\end{definition}

\section{Array}
\begin{definition}[array]~
    \begin{verbatim}
array=(v1 v2 ...)

for i in ${!array[@]}; do
    echo $array[$i]
done
\end{verbatim}
\end{definition}

\section{sed}

\section{awk}
\url{https://www.gnu.org/software/gawk/manual/gawk.html}

\begin{definition}[awk]
    The basic function of awk is to search files for lines that contain certain patterns.
\end{definition}

\begin{example}~
    \begin{verbatim}
# % @fig{png}
mapfile -t array < <(awk -F'[{}]' '/^% @fig\{/ { print $2 }' file)

# @book{BATCHELOR_2000_AN_INTRODUCTION_TO_FLUID_DYNAMICS,
awk -F'[{,]' '/@book\{/ {
    print "$1: " $1
    print "$2: " $2
    print "$3: " $3
    }' file
\end{verbatim}
\end{example}

\section{grep}
\begin{definition}[grep]
    Search input file for matches to the patterns.
\end{definition}

\section{find}

\section{Windows Batch}

\appendix
\section{Tools}
\small
\begin{verbatim}
#!/bin/bash

# Convert TikZ PDF to PNG
# pdf2png PDF PAGE PNG [DPI]
pdf2png() {
    local PDF=$1
    local PAGE=$2
    local PNG=$3
    local DPI=${4:-300}
    echo "Convert $PDF page $PAGE to $PNG"
    magick -density "$DPI" -units pixelsperinch \
        "$PDF[$PAGE]" \
        -alpha remove -background white -flatten \
        -trim +repage \
        -quality 100 \
        "$PNG"
}
\end{verbatim}

\begin{verbatim}
#!/bin/bash

# Compress PDF
# pdfcompress INPUT OUTPUT
pdfcompress() {
    local INPUT=$1
    local OUTPUT=$2
    ghostscript -sDEVICE=pdfwrite -dPDFSETTINGS=/prepress -dNOPAUSE -dQUIET -dBATCH \
        -sOUTPUTFile="$OUTPUT" "$INPUT"
}
\end{verbatim}

\printbibliography
\end{document}